% sec:geo-conservation — Conservation Monitoring
\section{Conservation Monitoring}
\label{sec:geo-conservation}

An adaptive integrator provides no guarantee that the trajectory it
returns is physically correct --- only that the local error at each step
satisfies the prescribed tolerance.  Accumulated discretisation error
can silently nudge a photon off its true geodesic path: a ray that
should graze the photon sphere and escape to the celestial sphere might
instead spiral inward, producing a shadow pixel where the accretion
disc should appear --- and nothing in the integrator flags the error.
Conservation laws provide an independent check.  Every quantity preserved by the exact equations
of motion but violated by the numerical approximation becomes a
diagnostic: the size of the violation measures the accumulated error
without requiring a reference solution.

\paragraph{The Hamiltonian diagnostic.}

The super-Hamiltonian $\mathcal{H} =
\tfrac{1}{2}\,\inversemetric\,p_\mu\,p_\nu$
(\cref{eq:geo-hamiltonian}) is exactly conserved along any geodesic
(\cref{eq:H-conservation}).  For null rays, $\mathcal{H} = 0$; for
massive particles, $\mathcal{H} = -\tfrac{1}{2}$
(\cref{eq:timelike-constraint}).  Any departure from this value is
pure discretisation error --- not a physical effect, not a coordinate
artefact, but a direct measure of how much the numerical trajectory
has deviated from a true geodesic.
\coderef{Geodesic.Hamiltonian}{hamiltonian}

The evaluation cost is negligible.  The inverse metric $\inversemetric$
is already computed at the current position for the position equation
\cref{eq:geo-xdot}, and the momentum $p_\mu$ is part of the state
vector.  The function \texttt{hamiltonian} contracts these to produce
the scalar $\mathcal{H}$:
%
\begin{lstlisting}[language=Haskell]
hamiltonian :: Sym4 'Contravariant -> Vec4 'Covariant -> Double
hamiltonian ginv mom =
  let pmUp = raise4 ginv mom
  in 0.5 * dot4 pmUp mom
\end{lstlisting}
%
One index raise and one inner product --- the same operations that
appear in the position equation, reused for monitoring.  The codebase
evaluates $\mathcal{H}$ after every accepted integration step, producing
a trace of constraint violation along the entire trajectory.
\implnote{The function \texttt{traceGeodesic} in
\codemodule{Geodesic.Solver} records triplets $(\lambda,
\text{state}, \mathcal{H})$ at every accepted step, providing the full
trajectory with Hamiltonian monitoring data for convergence studies and
visualisation.}

\paragraph{The error budget.}

The step-size controller (\cref{eq:rk45-stepcontrol}) keeps the local
truncation error at each step below a tolerance set by $\epsilon_a$
and $\epsilon_r$.  Over $N$ accepted steps, these local errors
accumulate in $\mathcal{H}$.  For a non-symplectic integrator, the
accumulation is not bounded a priori: it is a random walk when
successive local errors alternate in sign (total error
$\sim \sqrt{N}\,\epsilon$) and monotonic drift when the leading error
term maintains a consistent sign ($\sim N\,\epsilon$).  The
fifth-order accuracy of Dormand--Prince provides additional
suppression: the local truncation error of the propagated solution
scales as $\order{h^6}$, and over an integration interval of fixed
length the number of steps scales as $N \sim 1/h$, so the total
accumulated error scales as $N \cdot h^6 \sim h^5$.  In practice,
with the default tolerances
$\epsilon_a = \epsilon_r = 10^{-10}$, the Hamiltonian violation
for a typical deflected or plunging ray satisfies
%
\begin{equation}
  \abs{\mathcal{H}(\lambda) - \mathcal{H}_0}
  \lesssim \epsilon_a
  \label{eq:H-violation-bound}
\end{equation}
%
throughout the trajectory (\cref{fig:H-violation}) --- the violation
tracks the tolerance, as expected from the global error scaling.
Tightening the tolerance by two orders of magnitude reduces the
violation by the same factor, confirming the fifth-order convergence.
The tolerance can therefore be relaxed to $10^{-8}$ for preview
renders while guaranteeing $\abs{\mathcal{H}} \lesssim 10^{-8}$;
the $10^{-10}$ default ensures publication-quality accuracy.
\physnote{A symplectic integrator would produce bounded oscillations in
$\mathcal{H}$ rather than monotonic drift (backward error analysis).
The non-symplectic integrator trades this guarantee for adaptive step
control --- a favourable exchange for ray tracing, as argued in
\cref{sec:geo-raytracing}.}

\paragraph{Killing-vector quantities.}

The energy $E = -p_t$ and angular momentum $L = p_\varphi$ are conserved
along geodesics in any stationary, axisymmetric spacetime
(\cref{sec:geo-constants}).  Since $p_t$ and $p_\varphi$ are components
of the momentum covector --- already stored in the
\texttt{HamiltonianState} record --- monitoring them requires no extra
computation beyond reading state components.  For rays with nonzero
energy and angular momentum, the fractional drifts
%
\begin{equation}
  \frac{\abs{\Delta E}}{E_0}
  = \frac{\abs{p_t(\lambda) - p_t(0)}}{\abs{p_t(0)}} \,, \qquad
  \frac{\abs{\Delta L}}{L_0}
  = \frac{\abs{p_\varphi(\lambda) - p_\varphi(0)}}{\abs{p_\varphi(0)}}
  \label{eq:EL-drift}
\end{equation}
%
measure fractional energy and angular momentum error.  Typical values
for well-resolved ray traces satisfy
$\abs{\Delta E/E_0} \lesssim \epsilon_a$ and
$\abs{\Delta L/L_0} \lesssim \epsilon_a$, consistent with the
Hamiltonian bound \cref{eq:H-violation-bound}.

These diagnostics complement $\mathcal{H}$ because they probe
different aspects of the integration.  A coding error in the inverse
metric would be immediately visible in $\mathcal{H}$ (which depends
directly on $\inversemetric$), whereas its effect on $E$ and $L$ would
emerge only gradually through the integrated dynamics.  Conversely, a
spurious $\pd{t}\,g_{\alpha\beta}$ term (which should vanish in a
stationary spacetime) would produce $E$ drift with no corresponding
source in the exact equations, providing a more specific diagnostic
than the generic $\mathcal{H}$ violation.  Used together, the three
diagnostics cross-check both the mathematical implementation and the
numerical accuracy.
\implnote{In Kerr--Schild Cartesian coordinates, $E = -p_t$ is
read directly from the momentum covector since $t_{\text{KS}} =
t_{\text{BL}}$.  The angular momentum requires
$L = p_\varphi = x\,p_y - y\,p_x$, a chain-rule transformation
from Cartesian momenta.}

\paragraph{When monitoring signals trouble.}

Conservation violations grow fastest in two regimes.  Near the photon
sphere (\cref{sec:geo-raytracing}), the Lyapunov instability amplifies
local errors: the integrator shrinks the step size to maintain accuracy,
but each small error in position produces a disproportionate error in
the subsequent trajectory.  A near-critical ray with $b \approx b_c$
that winds several times around the photon sphere can accumulate
$\abs{\mathcal{H}} \sim 10\,\epsilon_a$ --- an order of magnitude above
the tolerance, and roughly ten times worse than a ray that passes
through the weak-field region without lingering
(\cref{fig:H-violation}).

The second regime is long integration intervals.  A ray traced from a
distant observer ($r_0 = 500M$) through a close approach at
$r \approx 3M$ and back out to $500M$ accumulates thousands of steps.
Even at $10^{-14}$ error per step, a thousand steps push the total to
$10^{-11}$.

\begin{figure}[t]
  \centering
  %% Creator: Matplotlib, PGF backend
%%
%% To include the figure in your LaTeX document, write
%%   \input{<filename>.pgf}
%%
%% Make sure the required packages are loaded in your preamble
%%   \usepackage{pgf}
%%
%% Also ensure that all the required font packages are loaded; for instance,
%% the lmodern package is sometimes necessary when using math font.
%%   \usepackage{lmodern}
%%
%% Figures using additional raster images can only be included by \input if
%% they are in the same directory as the main LaTeX file. For loading figures
%% from other directories you can use the `import` package
%%   \usepackage{import}
%%
%% and then include the figures with
%%   \import{<path to file>}{<filename>.pgf}
%%
%% Matplotlib used the following preamble
%%   \def\mathdefault#1{#1}
%%   \everymath=\expandafter{\the\everymath\displaystyle}
%%   \IfFileExists{scrextend.sty}{
%%     \usepackage[fontsize=10.000000pt]{scrextend}
%%   }{
%%     \renewcommand{\normalsize}{\fontsize{10.000000}{12.000000}\selectfont}
%%     \normalsize
%%   }
%%   \usepackage{fontspec}
%%   \usepackage{unicode-math}
%%   \setmainfont{STIX Two Text}[Numbers=OldStyle, Ligatures=TeX]
%%   \setmathfont{STIX Two Math}
%%   \setmonofont{Source Code Pro}[Scale=MatchLowercase]
%%   \ifdefined\pdftexversion\else  % non-pdftex case.
%%     \usepackage{fontspec}
%%     \setmainfont{DejaVuSerif.ttf}[Path=\detokenize{/Users/gvn/.pyenv/versions/3.12.11/lib/python3.12/site-packages/matplotlib/mpl-data/fonts/ttf/}]
%%     \setsansfont{DejaVuSans.ttf}[Path=\detokenize{/Users/gvn/.pyenv/versions/3.12.11/lib/python3.12/site-packages/matplotlib/mpl-data/fonts/ttf/}]
%%     \setmonofont{DejaVuSansMono.ttf}[Path=\detokenize{/Users/gvn/.pyenv/versions/3.12.11/lib/python3.12/site-packages/matplotlib/mpl-data/fonts/ttf/}]
%%   \fi
%%   \makeatletter\@ifpackageloaded{underscore}{}{\usepackage[strings]{underscore}}\makeatother
%%
\begingroup%
\makeatletter%
\begin{pgfpicture}%
\pgfpathrectangle{\pgfpointorigin}{\pgfqpoint{3.833993in}{1.991700in}}%
\pgfusepath{use as bounding box, clip}%
\begin{pgfscope}%
\pgfsetbuttcap%
\pgfsetmiterjoin%
\definecolor{currentfill}{rgb}{1.000000,1.000000,1.000000}%
\pgfsetfillcolor{currentfill}%
\pgfsetlinewidth{0.000000pt}%
\definecolor{currentstroke}{rgb}{1.000000,1.000000,1.000000}%
\pgfsetstrokecolor{currentstroke}%
\pgfsetdash{}{0pt}%
\pgfpathmoveto{\pgfqpoint{0.000000in}{0.000000in}}%
\pgfpathlineto{\pgfqpoint{3.833993in}{0.000000in}}%
\pgfpathlineto{\pgfqpoint{3.833993in}{1.991700in}}%
\pgfpathlineto{\pgfqpoint{0.000000in}{1.991700in}}%
\pgfpathlineto{\pgfqpoint{0.000000in}{0.000000in}}%
\pgfpathclose%
\pgfusepath{fill}%
\end{pgfscope}%
\begin{pgfscope}%
\pgfsetbuttcap%
\pgfsetmiterjoin%
\definecolor{currentfill}{rgb}{1.000000,1.000000,1.000000}%
\pgfsetfillcolor{currentfill}%
\pgfsetlinewidth{0.000000pt}%
\definecolor{currentstroke}{rgb}{0.000000,0.000000,0.000000}%
\pgfsetstrokecolor{currentstroke}%
\pgfsetstrokeopacity{0.000000}%
\pgfsetdash{}{0pt}%
\pgfpathmoveto{\pgfqpoint{0.530750in}{0.352669in}}%
\pgfpathlineto{\pgfqpoint{3.710853in}{0.352669in}}%
\pgfpathlineto{\pgfqpoint{3.710853in}{1.971700in}}%
\pgfpathlineto{\pgfqpoint{0.530750in}{1.971700in}}%
\pgfpathlineto{\pgfqpoint{0.530750in}{0.352669in}}%
\pgfpathclose%
\pgfusepath{fill}%
\end{pgfscope}%
\begin{pgfscope}%
\pgfpathrectangle{\pgfqpoint{0.530750in}{0.352669in}}{\pgfqpoint{3.180103in}{1.619031in}}%
\pgfusepath{clip}%
\pgfsetbuttcap%
\pgfsetroundjoin%
\pgfsetlinewidth{0.501875pt}%
\definecolor{currentstroke}{rgb}{0.431373,0.423529,0.462745}%
\pgfsetstrokecolor{currentstroke}%
\pgfsetdash{{1.850000pt}{0.800000pt}}{0.000000pt}%
\pgfpathmoveto{\pgfqpoint{0.530750in}{1.665888in}}%
\pgfpathlineto{\pgfqpoint{3.710853in}{1.665888in}}%
\pgfusepath{stroke}%
\end{pgfscope}%
\begin{pgfscope}%
\pgfsetbuttcap%
\pgfsetroundjoin%
\definecolor{currentfill}{rgb}{0.000000,0.000000,0.000000}%
\pgfsetfillcolor{currentfill}%
\pgfsetlinewidth{0.501875pt}%
\definecolor{currentstroke}{rgb}{0.000000,0.000000,0.000000}%
\pgfsetstrokecolor{currentstroke}%
\pgfsetdash{}{0pt}%
\pgfsys@defobject{currentmarker}{\pgfqpoint{0.000000in}{0.000000in}}{\pgfqpoint{0.000000in}{0.041667in}}{%
\pgfpathmoveto{\pgfqpoint{0.000000in}{0.000000in}}%
\pgfpathlineto{\pgfqpoint{0.000000in}{0.041667in}}%
\pgfusepath{stroke,fill}%
}%
\begin{pgfscope}%
\pgfsys@transformshift{0.675300in}{0.352669in}%
\pgfsys@useobject{currentmarker}{}%
\end{pgfscope}%
\end{pgfscope}%
\begin{pgfscope}%
\definecolor{textcolor}{rgb}{0.000000,0.000000,0.000000}%
\pgfsetstrokecolor{textcolor}%
\pgfsetfillcolor{textcolor}%
\pgftext[x=0.675300in,y=0.304058in,,top]{\color{textcolor}{\sffamily\fontsize{8.000000}{9.600000}\selectfont\catcode`\^=\active\def^{\ifmmode\sp\else\^{}\fi}\catcode`\%=\active\def%{\%}0}}%
\end{pgfscope}%
\begin{pgfscope}%
\pgfsetbuttcap%
\pgfsetroundjoin%
\definecolor{currentfill}{rgb}{0.000000,0.000000,0.000000}%
\pgfsetfillcolor{currentfill}%
\pgfsetlinewidth{0.501875pt}%
\definecolor{currentstroke}{rgb}{0.000000,0.000000,0.000000}%
\pgfsetstrokecolor{currentstroke}%
\pgfsetdash{}{0pt}%
\pgfsys@defobject{currentmarker}{\pgfqpoint{0.000000in}{0.000000in}}{\pgfqpoint{0.000000in}{0.041667in}}{%
\pgfpathmoveto{\pgfqpoint{0.000000in}{0.000000in}}%
\pgfpathlineto{\pgfqpoint{0.000000in}{0.041667in}}%
\pgfusepath{stroke,fill}%
}%
\begin{pgfscope}%
\pgfsys@transformshift{1.054382in}{0.352669in}%
\pgfsys@useobject{currentmarker}{}%
\end{pgfscope}%
\end{pgfscope}%
\begin{pgfscope}%
\definecolor{textcolor}{rgb}{0.000000,0.000000,0.000000}%
\pgfsetstrokecolor{textcolor}%
\pgfsetfillcolor{textcolor}%
\pgftext[x=1.054382in,y=0.304058in,,top]{\color{textcolor}{\sffamily\fontsize{8.000000}{9.600000}\selectfont\catcode`\^=\active\def^{\ifmmode\sp\else\^{}\fi}\catcode`\%=\active\def%{\%}20}}%
\end{pgfscope}%
\begin{pgfscope}%
\pgfsetbuttcap%
\pgfsetroundjoin%
\definecolor{currentfill}{rgb}{0.000000,0.000000,0.000000}%
\pgfsetfillcolor{currentfill}%
\pgfsetlinewidth{0.501875pt}%
\definecolor{currentstroke}{rgb}{0.000000,0.000000,0.000000}%
\pgfsetstrokecolor{currentstroke}%
\pgfsetdash{}{0pt}%
\pgfsys@defobject{currentmarker}{\pgfqpoint{0.000000in}{0.000000in}}{\pgfqpoint{0.000000in}{0.041667in}}{%
\pgfpathmoveto{\pgfqpoint{0.000000in}{0.000000in}}%
\pgfpathlineto{\pgfqpoint{0.000000in}{0.041667in}}%
\pgfusepath{stroke,fill}%
}%
\begin{pgfscope}%
\pgfsys@transformshift{1.433464in}{0.352669in}%
\pgfsys@useobject{currentmarker}{}%
\end{pgfscope}%
\end{pgfscope}%
\begin{pgfscope}%
\definecolor{textcolor}{rgb}{0.000000,0.000000,0.000000}%
\pgfsetstrokecolor{textcolor}%
\pgfsetfillcolor{textcolor}%
\pgftext[x=1.433464in,y=0.304058in,,top]{\color{textcolor}{\sffamily\fontsize{8.000000}{9.600000}\selectfont\catcode`\^=\active\def^{\ifmmode\sp\else\^{}\fi}\catcode`\%=\active\def%{\%}40}}%
\end{pgfscope}%
\begin{pgfscope}%
\pgfsetbuttcap%
\pgfsetroundjoin%
\definecolor{currentfill}{rgb}{0.000000,0.000000,0.000000}%
\pgfsetfillcolor{currentfill}%
\pgfsetlinewidth{0.501875pt}%
\definecolor{currentstroke}{rgb}{0.000000,0.000000,0.000000}%
\pgfsetstrokecolor{currentstroke}%
\pgfsetdash{}{0pt}%
\pgfsys@defobject{currentmarker}{\pgfqpoint{0.000000in}{0.000000in}}{\pgfqpoint{0.000000in}{0.041667in}}{%
\pgfpathmoveto{\pgfqpoint{0.000000in}{0.000000in}}%
\pgfpathlineto{\pgfqpoint{0.000000in}{0.041667in}}%
\pgfusepath{stroke,fill}%
}%
\begin{pgfscope}%
\pgfsys@transformshift{1.812545in}{0.352669in}%
\pgfsys@useobject{currentmarker}{}%
\end{pgfscope}%
\end{pgfscope}%
\begin{pgfscope}%
\definecolor{textcolor}{rgb}{0.000000,0.000000,0.000000}%
\pgfsetstrokecolor{textcolor}%
\pgfsetfillcolor{textcolor}%
\pgftext[x=1.812545in,y=0.304058in,,top]{\color{textcolor}{\sffamily\fontsize{8.000000}{9.600000}\selectfont\catcode`\^=\active\def^{\ifmmode\sp\else\^{}\fi}\catcode`\%=\active\def%{\%}60}}%
\end{pgfscope}%
\begin{pgfscope}%
\pgfsetbuttcap%
\pgfsetroundjoin%
\definecolor{currentfill}{rgb}{0.000000,0.000000,0.000000}%
\pgfsetfillcolor{currentfill}%
\pgfsetlinewidth{0.501875pt}%
\definecolor{currentstroke}{rgb}{0.000000,0.000000,0.000000}%
\pgfsetstrokecolor{currentstroke}%
\pgfsetdash{}{0pt}%
\pgfsys@defobject{currentmarker}{\pgfqpoint{0.000000in}{0.000000in}}{\pgfqpoint{0.000000in}{0.041667in}}{%
\pgfpathmoveto{\pgfqpoint{0.000000in}{0.000000in}}%
\pgfpathlineto{\pgfqpoint{0.000000in}{0.041667in}}%
\pgfusepath{stroke,fill}%
}%
\begin{pgfscope}%
\pgfsys@transformshift{2.191627in}{0.352669in}%
\pgfsys@useobject{currentmarker}{}%
\end{pgfscope}%
\end{pgfscope}%
\begin{pgfscope}%
\definecolor{textcolor}{rgb}{0.000000,0.000000,0.000000}%
\pgfsetstrokecolor{textcolor}%
\pgfsetfillcolor{textcolor}%
\pgftext[x=2.191627in,y=0.304058in,,top]{\color{textcolor}{\sffamily\fontsize{8.000000}{9.600000}\selectfont\catcode`\^=\active\def^{\ifmmode\sp\else\^{}\fi}\catcode`\%=\active\def%{\%}80}}%
\end{pgfscope}%
\begin{pgfscope}%
\pgfsetbuttcap%
\pgfsetroundjoin%
\definecolor{currentfill}{rgb}{0.000000,0.000000,0.000000}%
\pgfsetfillcolor{currentfill}%
\pgfsetlinewidth{0.501875pt}%
\definecolor{currentstroke}{rgb}{0.000000,0.000000,0.000000}%
\pgfsetstrokecolor{currentstroke}%
\pgfsetdash{}{0pt}%
\pgfsys@defobject{currentmarker}{\pgfqpoint{0.000000in}{0.000000in}}{\pgfqpoint{0.000000in}{0.041667in}}{%
\pgfpathmoveto{\pgfqpoint{0.000000in}{0.000000in}}%
\pgfpathlineto{\pgfqpoint{0.000000in}{0.041667in}}%
\pgfusepath{stroke,fill}%
}%
\begin{pgfscope}%
\pgfsys@transformshift{2.570709in}{0.352669in}%
\pgfsys@useobject{currentmarker}{}%
\end{pgfscope}%
\end{pgfscope}%
\begin{pgfscope}%
\definecolor{textcolor}{rgb}{0.000000,0.000000,0.000000}%
\pgfsetstrokecolor{textcolor}%
\pgfsetfillcolor{textcolor}%
\pgftext[x=2.570709in,y=0.304058in,,top]{\color{textcolor}{\sffamily\fontsize{8.000000}{9.600000}\selectfont\catcode`\^=\active\def^{\ifmmode\sp\else\^{}\fi}\catcode`\%=\active\def%{\%}100}}%
\end{pgfscope}%
\begin{pgfscope}%
\pgfsetbuttcap%
\pgfsetroundjoin%
\definecolor{currentfill}{rgb}{0.000000,0.000000,0.000000}%
\pgfsetfillcolor{currentfill}%
\pgfsetlinewidth{0.501875pt}%
\definecolor{currentstroke}{rgb}{0.000000,0.000000,0.000000}%
\pgfsetstrokecolor{currentstroke}%
\pgfsetdash{}{0pt}%
\pgfsys@defobject{currentmarker}{\pgfqpoint{0.000000in}{0.000000in}}{\pgfqpoint{0.000000in}{0.041667in}}{%
\pgfpathmoveto{\pgfqpoint{0.000000in}{0.000000in}}%
\pgfpathlineto{\pgfqpoint{0.000000in}{0.041667in}}%
\pgfusepath{stroke,fill}%
}%
\begin{pgfscope}%
\pgfsys@transformshift{2.949791in}{0.352669in}%
\pgfsys@useobject{currentmarker}{}%
\end{pgfscope}%
\end{pgfscope}%
\begin{pgfscope}%
\definecolor{textcolor}{rgb}{0.000000,0.000000,0.000000}%
\pgfsetstrokecolor{textcolor}%
\pgfsetfillcolor{textcolor}%
\pgftext[x=2.949791in,y=0.304058in,,top]{\color{textcolor}{\sffamily\fontsize{8.000000}{9.600000}\selectfont\catcode`\^=\active\def^{\ifmmode\sp\else\^{}\fi}\catcode`\%=\active\def%{\%}120}}%
\end{pgfscope}%
\begin{pgfscope}%
\pgfsetbuttcap%
\pgfsetroundjoin%
\definecolor{currentfill}{rgb}{0.000000,0.000000,0.000000}%
\pgfsetfillcolor{currentfill}%
\pgfsetlinewidth{0.501875pt}%
\definecolor{currentstroke}{rgb}{0.000000,0.000000,0.000000}%
\pgfsetstrokecolor{currentstroke}%
\pgfsetdash{}{0pt}%
\pgfsys@defobject{currentmarker}{\pgfqpoint{0.000000in}{0.000000in}}{\pgfqpoint{0.000000in}{0.041667in}}{%
\pgfpathmoveto{\pgfqpoint{0.000000in}{0.000000in}}%
\pgfpathlineto{\pgfqpoint{0.000000in}{0.041667in}}%
\pgfusepath{stroke,fill}%
}%
\begin{pgfscope}%
\pgfsys@transformshift{3.328873in}{0.352669in}%
\pgfsys@useobject{currentmarker}{}%
\end{pgfscope}%
\end{pgfscope}%
\begin{pgfscope}%
\definecolor{textcolor}{rgb}{0.000000,0.000000,0.000000}%
\pgfsetstrokecolor{textcolor}%
\pgfsetfillcolor{textcolor}%
\pgftext[x=3.328873in,y=0.304058in,,top]{\color{textcolor}{\sffamily\fontsize{8.000000}{9.600000}\selectfont\catcode`\^=\active\def^{\ifmmode\sp\else\^{}\fi}\catcode`\%=\active\def%{\%}140}}%
\end{pgfscope}%
\begin{pgfscope}%
\pgfsetbuttcap%
\pgfsetroundjoin%
\definecolor{currentfill}{rgb}{0.000000,0.000000,0.000000}%
\pgfsetfillcolor{currentfill}%
\pgfsetlinewidth{0.501875pt}%
\definecolor{currentstroke}{rgb}{0.000000,0.000000,0.000000}%
\pgfsetstrokecolor{currentstroke}%
\pgfsetdash{}{0pt}%
\pgfsys@defobject{currentmarker}{\pgfqpoint{0.000000in}{0.000000in}}{\pgfqpoint{0.000000in}{0.041667in}}{%
\pgfpathmoveto{\pgfqpoint{0.000000in}{0.000000in}}%
\pgfpathlineto{\pgfqpoint{0.000000in}{0.041667in}}%
\pgfusepath{stroke,fill}%
}%
\begin{pgfscope}%
\pgfsys@transformshift{3.707955in}{0.352669in}%
\pgfsys@useobject{currentmarker}{}%
\end{pgfscope}%
\end{pgfscope}%
\begin{pgfscope}%
\definecolor{textcolor}{rgb}{0.000000,0.000000,0.000000}%
\pgfsetstrokecolor{textcolor}%
\pgfsetfillcolor{textcolor}%
\pgftext[x=3.707955in,y=0.304058in,,top]{\color{textcolor}{\sffamily\fontsize{8.000000}{9.600000}\selectfont\catcode`\^=\active\def^{\ifmmode\sp\else\^{}\fi}\catcode`\%=\active\def%{\%}160}}%
\end{pgfscope}%
\begin{pgfscope}%
\pgfsetbuttcap%
\pgfsetroundjoin%
\definecolor{currentfill}{rgb}{0.000000,0.000000,0.000000}%
\pgfsetfillcolor{currentfill}%
\pgfsetlinewidth{0.301125pt}%
\definecolor{currentstroke}{rgb}{0.000000,0.000000,0.000000}%
\pgfsetstrokecolor{currentstroke}%
\pgfsetdash{}{0pt}%
\pgfsys@defobject{currentmarker}{\pgfqpoint{0.000000in}{0.000000in}}{\pgfqpoint{0.000000in}{0.027778in}}{%
\pgfpathmoveto{\pgfqpoint{0.000000in}{0.000000in}}%
\pgfpathlineto{\pgfqpoint{0.000000in}{0.027778in}}%
\pgfusepath{stroke,fill}%
}%
\begin{pgfscope}%
\pgfsys@transformshift{0.580529in}{0.352669in}%
\pgfsys@useobject{currentmarker}{}%
\end{pgfscope}%
\end{pgfscope}%
\begin{pgfscope}%
\pgfsetbuttcap%
\pgfsetroundjoin%
\definecolor{currentfill}{rgb}{0.000000,0.000000,0.000000}%
\pgfsetfillcolor{currentfill}%
\pgfsetlinewidth{0.301125pt}%
\definecolor{currentstroke}{rgb}{0.000000,0.000000,0.000000}%
\pgfsetstrokecolor{currentstroke}%
\pgfsetdash{}{0pt}%
\pgfsys@defobject{currentmarker}{\pgfqpoint{0.000000in}{0.000000in}}{\pgfqpoint{0.000000in}{0.027778in}}{%
\pgfpathmoveto{\pgfqpoint{0.000000in}{0.000000in}}%
\pgfpathlineto{\pgfqpoint{0.000000in}{0.027778in}}%
\pgfusepath{stroke,fill}%
}%
\begin{pgfscope}%
\pgfsys@transformshift{0.770070in}{0.352669in}%
\pgfsys@useobject{currentmarker}{}%
\end{pgfscope}%
\end{pgfscope}%
\begin{pgfscope}%
\pgfsetbuttcap%
\pgfsetroundjoin%
\definecolor{currentfill}{rgb}{0.000000,0.000000,0.000000}%
\pgfsetfillcolor{currentfill}%
\pgfsetlinewidth{0.301125pt}%
\definecolor{currentstroke}{rgb}{0.000000,0.000000,0.000000}%
\pgfsetstrokecolor{currentstroke}%
\pgfsetdash{}{0pt}%
\pgfsys@defobject{currentmarker}{\pgfqpoint{0.000000in}{0.000000in}}{\pgfqpoint{0.000000in}{0.027778in}}{%
\pgfpathmoveto{\pgfqpoint{0.000000in}{0.000000in}}%
\pgfpathlineto{\pgfqpoint{0.000000in}{0.027778in}}%
\pgfusepath{stroke,fill}%
}%
\begin{pgfscope}%
\pgfsys@transformshift{0.864841in}{0.352669in}%
\pgfsys@useobject{currentmarker}{}%
\end{pgfscope}%
\end{pgfscope}%
\begin{pgfscope}%
\pgfsetbuttcap%
\pgfsetroundjoin%
\definecolor{currentfill}{rgb}{0.000000,0.000000,0.000000}%
\pgfsetfillcolor{currentfill}%
\pgfsetlinewidth{0.301125pt}%
\definecolor{currentstroke}{rgb}{0.000000,0.000000,0.000000}%
\pgfsetstrokecolor{currentstroke}%
\pgfsetdash{}{0pt}%
\pgfsys@defobject{currentmarker}{\pgfqpoint{0.000000in}{0.000000in}}{\pgfqpoint{0.000000in}{0.027778in}}{%
\pgfpathmoveto{\pgfqpoint{0.000000in}{0.000000in}}%
\pgfpathlineto{\pgfqpoint{0.000000in}{0.027778in}}%
\pgfusepath{stroke,fill}%
}%
\begin{pgfscope}%
\pgfsys@transformshift{0.959611in}{0.352669in}%
\pgfsys@useobject{currentmarker}{}%
\end{pgfscope}%
\end{pgfscope}%
\begin{pgfscope}%
\pgfsetbuttcap%
\pgfsetroundjoin%
\definecolor{currentfill}{rgb}{0.000000,0.000000,0.000000}%
\pgfsetfillcolor{currentfill}%
\pgfsetlinewidth{0.301125pt}%
\definecolor{currentstroke}{rgb}{0.000000,0.000000,0.000000}%
\pgfsetstrokecolor{currentstroke}%
\pgfsetdash{}{0pt}%
\pgfsys@defobject{currentmarker}{\pgfqpoint{0.000000in}{0.000000in}}{\pgfqpoint{0.000000in}{0.027778in}}{%
\pgfpathmoveto{\pgfqpoint{0.000000in}{0.000000in}}%
\pgfpathlineto{\pgfqpoint{0.000000in}{0.027778in}}%
\pgfusepath{stroke,fill}%
}%
\begin{pgfscope}%
\pgfsys@transformshift{1.149152in}{0.352669in}%
\pgfsys@useobject{currentmarker}{}%
\end{pgfscope}%
\end{pgfscope}%
\begin{pgfscope}%
\pgfsetbuttcap%
\pgfsetroundjoin%
\definecolor{currentfill}{rgb}{0.000000,0.000000,0.000000}%
\pgfsetfillcolor{currentfill}%
\pgfsetlinewidth{0.301125pt}%
\definecolor{currentstroke}{rgb}{0.000000,0.000000,0.000000}%
\pgfsetstrokecolor{currentstroke}%
\pgfsetdash{}{0pt}%
\pgfsys@defobject{currentmarker}{\pgfqpoint{0.000000in}{0.000000in}}{\pgfqpoint{0.000000in}{0.027778in}}{%
\pgfpathmoveto{\pgfqpoint{0.000000in}{0.000000in}}%
\pgfpathlineto{\pgfqpoint{0.000000in}{0.027778in}}%
\pgfusepath{stroke,fill}%
}%
\begin{pgfscope}%
\pgfsys@transformshift{1.243923in}{0.352669in}%
\pgfsys@useobject{currentmarker}{}%
\end{pgfscope}%
\end{pgfscope}%
\begin{pgfscope}%
\pgfsetbuttcap%
\pgfsetroundjoin%
\definecolor{currentfill}{rgb}{0.000000,0.000000,0.000000}%
\pgfsetfillcolor{currentfill}%
\pgfsetlinewidth{0.301125pt}%
\definecolor{currentstroke}{rgb}{0.000000,0.000000,0.000000}%
\pgfsetstrokecolor{currentstroke}%
\pgfsetdash{}{0pt}%
\pgfsys@defobject{currentmarker}{\pgfqpoint{0.000000in}{0.000000in}}{\pgfqpoint{0.000000in}{0.027778in}}{%
\pgfpathmoveto{\pgfqpoint{0.000000in}{0.000000in}}%
\pgfpathlineto{\pgfqpoint{0.000000in}{0.027778in}}%
\pgfusepath{stroke,fill}%
}%
\begin{pgfscope}%
\pgfsys@transformshift{1.338693in}{0.352669in}%
\pgfsys@useobject{currentmarker}{}%
\end{pgfscope}%
\end{pgfscope}%
\begin{pgfscope}%
\pgfsetbuttcap%
\pgfsetroundjoin%
\definecolor{currentfill}{rgb}{0.000000,0.000000,0.000000}%
\pgfsetfillcolor{currentfill}%
\pgfsetlinewidth{0.301125pt}%
\definecolor{currentstroke}{rgb}{0.000000,0.000000,0.000000}%
\pgfsetstrokecolor{currentstroke}%
\pgfsetdash{}{0pt}%
\pgfsys@defobject{currentmarker}{\pgfqpoint{0.000000in}{0.000000in}}{\pgfqpoint{0.000000in}{0.027778in}}{%
\pgfpathmoveto{\pgfqpoint{0.000000in}{0.000000in}}%
\pgfpathlineto{\pgfqpoint{0.000000in}{0.027778in}}%
\pgfusepath{stroke,fill}%
}%
\begin{pgfscope}%
\pgfsys@transformshift{1.528234in}{0.352669in}%
\pgfsys@useobject{currentmarker}{}%
\end{pgfscope}%
\end{pgfscope}%
\begin{pgfscope}%
\pgfsetbuttcap%
\pgfsetroundjoin%
\definecolor{currentfill}{rgb}{0.000000,0.000000,0.000000}%
\pgfsetfillcolor{currentfill}%
\pgfsetlinewidth{0.301125pt}%
\definecolor{currentstroke}{rgb}{0.000000,0.000000,0.000000}%
\pgfsetstrokecolor{currentstroke}%
\pgfsetdash{}{0pt}%
\pgfsys@defobject{currentmarker}{\pgfqpoint{0.000000in}{0.000000in}}{\pgfqpoint{0.000000in}{0.027778in}}{%
\pgfpathmoveto{\pgfqpoint{0.000000in}{0.000000in}}%
\pgfpathlineto{\pgfqpoint{0.000000in}{0.027778in}}%
\pgfusepath{stroke,fill}%
}%
\begin{pgfscope}%
\pgfsys@transformshift{1.623004in}{0.352669in}%
\pgfsys@useobject{currentmarker}{}%
\end{pgfscope}%
\end{pgfscope}%
\begin{pgfscope}%
\pgfsetbuttcap%
\pgfsetroundjoin%
\definecolor{currentfill}{rgb}{0.000000,0.000000,0.000000}%
\pgfsetfillcolor{currentfill}%
\pgfsetlinewidth{0.301125pt}%
\definecolor{currentstroke}{rgb}{0.000000,0.000000,0.000000}%
\pgfsetstrokecolor{currentstroke}%
\pgfsetdash{}{0pt}%
\pgfsys@defobject{currentmarker}{\pgfqpoint{0.000000in}{0.000000in}}{\pgfqpoint{0.000000in}{0.027778in}}{%
\pgfpathmoveto{\pgfqpoint{0.000000in}{0.000000in}}%
\pgfpathlineto{\pgfqpoint{0.000000in}{0.027778in}}%
\pgfusepath{stroke,fill}%
}%
\begin{pgfscope}%
\pgfsys@transformshift{1.717775in}{0.352669in}%
\pgfsys@useobject{currentmarker}{}%
\end{pgfscope}%
\end{pgfscope}%
\begin{pgfscope}%
\pgfsetbuttcap%
\pgfsetroundjoin%
\definecolor{currentfill}{rgb}{0.000000,0.000000,0.000000}%
\pgfsetfillcolor{currentfill}%
\pgfsetlinewidth{0.301125pt}%
\definecolor{currentstroke}{rgb}{0.000000,0.000000,0.000000}%
\pgfsetstrokecolor{currentstroke}%
\pgfsetdash{}{0pt}%
\pgfsys@defobject{currentmarker}{\pgfqpoint{0.000000in}{0.000000in}}{\pgfqpoint{0.000000in}{0.027778in}}{%
\pgfpathmoveto{\pgfqpoint{0.000000in}{0.000000in}}%
\pgfpathlineto{\pgfqpoint{0.000000in}{0.027778in}}%
\pgfusepath{stroke,fill}%
}%
\begin{pgfscope}%
\pgfsys@transformshift{1.907316in}{0.352669in}%
\pgfsys@useobject{currentmarker}{}%
\end{pgfscope}%
\end{pgfscope}%
\begin{pgfscope}%
\pgfsetbuttcap%
\pgfsetroundjoin%
\definecolor{currentfill}{rgb}{0.000000,0.000000,0.000000}%
\pgfsetfillcolor{currentfill}%
\pgfsetlinewidth{0.301125pt}%
\definecolor{currentstroke}{rgb}{0.000000,0.000000,0.000000}%
\pgfsetstrokecolor{currentstroke}%
\pgfsetdash{}{0pt}%
\pgfsys@defobject{currentmarker}{\pgfqpoint{0.000000in}{0.000000in}}{\pgfqpoint{0.000000in}{0.027778in}}{%
\pgfpathmoveto{\pgfqpoint{0.000000in}{0.000000in}}%
\pgfpathlineto{\pgfqpoint{0.000000in}{0.027778in}}%
\pgfusepath{stroke,fill}%
}%
\begin{pgfscope}%
\pgfsys@transformshift{2.002086in}{0.352669in}%
\pgfsys@useobject{currentmarker}{}%
\end{pgfscope}%
\end{pgfscope}%
\begin{pgfscope}%
\pgfsetbuttcap%
\pgfsetroundjoin%
\definecolor{currentfill}{rgb}{0.000000,0.000000,0.000000}%
\pgfsetfillcolor{currentfill}%
\pgfsetlinewidth{0.301125pt}%
\definecolor{currentstroke}{rgb}{0.000000,0.000000,0.000000}%
\pgfsetstrokecolor{currentstroke}%
\pgfsetdash{}{0pt}%
\pgfsys@defobject{currentmarker}{\pgfqpoint{0.000000in}{0.000000in}}{\pgfqpoint{0.000000in}{0.027778in}}{%
\pgfpathmoveto{\pgfqpoint{0.000000in}{0.000000in}}%
\pgfpathlineto{\pgfqpoint{0.000000in}{0.027778in}}%
\pgfusepath{stroke,fill}%
}%
\begin{pgfscope}%
\pgfsys@transformshift{2.096857in}{0.352669in}%
\pgfsys@useobject{currentmarker}{}%
\end{pgfscope}%
\end{pgfscope}%
\begin{pgfscope}%
\pgfsetbuttcap%
\pgfsetroundjoin%
\definecolor{currentfill}{rgb}{0.000000,0.000000,0.000000}%
\pgfsetfillcolor{currentfill}%
\pgfsetlinewidth{0.301125pt}%
\definecolor{currentstroke}{rgb}{0.000000,0.000000,0.000000}%
\pgfsetstrokecolor{currentstroke}%
\pgfsetdash{}{0pt}%
\pgfsys@defobject{currentmarker}{\pgfqpoint{0.000000in}{0.000000in}}{\pgfqpoint{0.000000in}{0.027778in}}{%
\pgfpathmoveto{\pgfqpoint{0.000000in}{0.000000in}}%
\pgfpathlineto{\pgfqpoint{0.000000in}{0.027778in}}%
\pgfusepath{stroke,fill}%
}%
\begin{pgfscope}%
\pgfsys@transformshift{2.286398in}{0.352669in}%
\pgfsys@useobject{currentmarker}{}%
\end{pgfscope}%
\end{pgfscope}%
\begin{pgfscope}%
\pgfsetbuttcap%
\pgfsetroundjoin%
\definecolor{currentfill}{rgb}{0.000000,0.000000,0.000000}%
\pgfsetfillcolor{currentfill}%
\pgfsetlinewidth{0.301125pt}%
\definecolor{currentstroke}{rgb}{0.000000,0.000000,0.000000}%
\pgfsetstrokecolor{currentstroke}%
\pgfsetdash{}{0pt}%
\pgfsys@defobject{currentmarker}{\pgfqpoint{0.000000in}{0.000000in}}{\pgfqpoint{0.000000in}{0.027778in}}{%
\pgfpathmoveto{\pgfqpoint{0.000000in}{0.000000in}}%
\pgfpathlineto{\pgfqpoint{0.000000in}{0.027778in}}%
\pgfusepath{stroke,fill}%
}%
\begin{pgfscope}%
\pgfsys@transformshift{2.381168in}{0.352669in}%
\pgfsys@useobject{currentmarker}{}%
\end{pgfscope}%
\end{pgfscope}%
\begin{pgfscope}%
\pgfsetbuttcap%
\pgfsetroundjoin%
\definecolor{currentfill}{rgb}{0.000000,0.000000,0.000000}%
\pgfsetfillcolor{currentfill}%
\pgfsetlinewidth{0.301125pt}%
\definecolor{currentstroke}{rgb}{0.000000,0.000000,0.000000}%
\pgfsetstrokecolor{currentstroke}%
\pgfsetdash{}{0pt}%
\pgfsys@defobject{currentmarker}{\pgfqpoint{0.000000in}{0.000000in}}{\pgfqpoint{0.000000in}{0.027778in}}{%
\pgfpathmoveto{\pgfqpoint{0.000000in}{0.000000in}}%
\pgfpathlineto{\pgfqpoint{0.000000in}{0.027778in}}%
\pgfusepath{stroke,fill}%
}%
\begin{pgfscope}%
\pgfsys@transformshift{2.475939in}{0.352669in}%
\pgfsys@useobject{currentmarker}{}%
\end{pgfscope}%
\end{pgfscope}%
\begin{pgfscope}%
\pgfsetbuttcap%
\pgfsetroundjoin%
\definecolor{currentfill}{rgb}{0.000000,0.000000,0.000000}%
\pgfsetfillcolor{currentfill}%
\pgfsetlinewidth{0.301125pt}%
\definecolor{currentstroke}{rgb}{0.000000,0.000000,0.000000}%
\pgfsetstrokecolor{currentstroke}%
\pgfsetdash{}{0pt}%
\pgfsys@defobject{currentmarker}{\pgfqpoint{0.000000in}{0.000000in}}{\pgfqpoint{0.000000in}{0.027778in}}{%
\pgfpathmoveto{\pgfqpoint{0.000000in}{0.000000in}}%
\pgfpathlineto{\pgfqpoint{0.000000in}{0.027778in}}%
\pgfusepath{stroke,fill}%
}%
\begin{pgfscope}%
\pgfsys@transformshift{2.665480in}{0.352669in}%
\pgfsys@useobject{currentmarker}{}%
\end{pgfscope}%
\end{pgfscope}%
\begin{pgfscope}%
\pgfsetbuttcap%
\pgfsetroundjoin%
\definecolor{currentfill}{rgb}{0.000000,0.000000,0.000000}%
\pgfsetfillcolor{currentfill}%
\pgfsetlinewidth{0.301125pt}%
\definecolor{currentstroke}{rgb}{0.000000,0.000000,0.000000}%
\pgfsetstrokecolor{currentstroke}%
\pgfsetdash{}{0pt}%
\pgfsys@defobject{currentmarker}{\pgfqpoint{0.000000in}{0.000000in}}{\pgfqpoint{0.000000in}{0.027778in}}{%
\pgfpathmoveto{\pgfqpoint{0.000000in}{0.000000in}}%
\pgfpathlineto{\pgfqpoint{0.000000in}{0.027778in}}%
\pgfusepath{stroke,fill}%
}%
\begin{pgfscope}%
\pgfsys@transformshift{2.760250in}{0.352669in}%
\pgfsys@useobject{currentmarker}{}%
\end{pgfscope}%
\end{pgfscope}%
\begin{pgfscope}%
\pgfsetbuttcap%
\pgfsetroundjoin%
\definecolor{currentfill}{rgb}{0.000000,0.000000,0.000000}%
\pgfsetfillcolor{currentfill}%
\pgfsetlinewidth{0.301125pt}%
\definecolor{currentstroke}{rgb}{0.000000,0.000000,0.000000}%
\pgfsetstrokecolor{currentstroke}%
\pgfsetdash{}{0pt}%
\pgfsys@defobject{currentmarker}{\pgfqpoint{0.000000in}{0.000000in}}{\pgfqpoint{0.000000in}{0.027778in}}{%
\pgfpathmoveto{\pgfqpoint{0.000000in}{0.000000in}}%
\pgfpathlineto{\pgfqpoint{0.000000in}{0.027778in}}%
\pgfusepath{stroke,fill}%
}%
\begin{pgfscope}%
\pgfsys@transformshift{2.855021in}{0.352669in}%
\pgfsys@useobject{currentmarker}{}%
\end{pgfscope}%
\end{pgfscope}%
\begin{pgfscope}%
\pgfsetbuttcap%
\pgfsetroundjoin%
\definecolor{currentfill}{rgb}{0.000000,0.000000,0.000000}%
\pgfsetfillcolor{currentfill}%
\pgfsetlinewidth{0.301125pt}%
\definecolor{currentstroke}{rgb}{0.000000,0.000000,0.000000}%
\pgfsetstrokecolor{currentstroke}%
\pgfsetdash{}{0pt}%
\pgfsys@defobject{currentmarker}{\pgfqpoint{0.000000in}{0.000000in}}{\pgfqpoint{0.000000in}{0.027778in}}{%
\pgfpathmoveto{\pgfqpoint{0.000000in}{0.000000in}}%
\pgfpathlineto{\pgfqpoint{0.000000in}{0.027778in}}%
\pgfusepath{stroke,fill}%
}%
\begin{pgfscope}%
\pgfsys@transformshift{3.044562in}{0.352669in}%
\pgfsys@useobject{currentmarker}{}%
\end{pgfscope}%
\end{pgfscope}%
\begin{pgfscope}%
\pgfsetbuttcap%
\pgfsetroundjoin%
\definecolor{currentfill}{rgb}{0.000000,0.000000,0.000000}%
\pgfsetfillcolor{currentfill}%
\pgfsetlinewidth{0.301125pt}%
\definecolor{currentstroke}{rgb}{0.000000,0.000000,0.000000}%
\pgfsetstrokecolor{currentstroke}%
\pgfsetdash{}{0pt}%
\pgfsys@defobject{currentmarker}{\pgfqpoint{0.000000in}{0.000000in}}{\pgfqpoint{0.000000in}{0.027778in}}{%
\pgfpathmoveto{\pgfqpoint{0.000000in}{0.000000in}}%
\pgfpathlineto{\pgfqpoint{0.000000in}{0.027778in}}%
\pgfusepath{stroke,fill}%
}%
\begin{pgfscope}%
\pgfsys@transformshift{3.139332in}{0.352669in}%
\pgfsys@useobject{currentmarker}{}%
\end{pgfscope}%
\end{pgfscope}%
\begin{pgfscope}%
\pgfsetbuttcap%
\pgfsetroundjoin%
\definecolor{currentfill}{rgb}{0.000000,0.000000,0.000000}%
\pgfsetfillcolor{currentfill}%
\pgfsetlinewidth{0.301125pt}%
\definecolor{currentstroke}{rgb}{0.000000,0.000000,0.000000}%
\pgfsetstrokecolor{currentstroke}%
\pgfsetdash{}{0pt}%
\pgfsys@defobject{currentmarker}{\pgfqpoint{0.000000in}{0.000000in}}{\pgfqpoint{0.000000in}{0.027778in}}{%
\pgfpathmoveto{\pgfqpoint{0.000000in}{0.000000in}}%
\pgfpathlineto{\pgfqpoint{0.000000in}{0.027778in}}%
\pgfusepath{stroke,fill}%
}%
\begin{pgfscope}%
\pgfsys@transformshift{3.234103in}{0.352669in}%
\pgfsys@useobject{currentmarker}{}%
\end{pgfscope}%
\end{pgfscope}%
\begin{pgfscope}%
\pgfsetbuttcap%
\pgfsetroundjoin%
\definecolor{currentfill}{rgb}{0.000000,0.000000,0.000000}%
\pgfsetfillcolor{currentfill}%
\pgfsetlinewidth{0.301125pt}%
\definecolor{currentstroke}{rgb}{0.000000,0.000000,0.000000}%
\pgfsetstrokecolor{currentstroke}%
\pgfsetdash{}{0pt}%
\pgfsys@defobject{currentmarker}{\pgfqpoint{0.000000in}{0.000000in}}{\pgfqpoint{0.000000in}{0.027778in}}{%
\pgfpathmoveto{\pgfqpoint{0.000000in}{0.000000in}}%
\pgfpathlineto{\pgfqpoint{0.000000in}{0.027778in}}%
\pgfusepath{stroke,fill}%
}%
\begin{pgfscope}%
\pgfsys@transformshift{3.423644in}{0.352669in}%
\pgfsys@useobject{currentmarker}{}%
\end{pgfscope}%
\end{pgfscope}%
\begin{pgfscope}%
\pgfsetbuttcap%
\pgfsetroundjoin%
\definecolor{currentfill}{rgb}{0.000000,0.000000,0.000000}%
\pgfsetfillcolor{currentfill}%
\pgfsetlinewidth{0.301125pt}%
\definecolor{currentstroke}{rgb}{0.000000,0.000000,0.000000}%
\pgfsetstrokecolor{currentstroke}%
\pgfsetdash{}{0pt}%
\pgfsys@defobject{currentmarker}{\pgfqpoint{0.000000in}{0.000000in}}{\pgfqpoint{0.000000in}{0.027778in}}{%
\pgfpathmoveto{\pgfqpoint{0.000000in}{0.000000in}}%
\pgfpathlineto{\pgfqpoint{0.000000in}{0.027778in}}%
\pgfusepath{stroke,fill}%
}%
\begin{pgfscope}%
\pgfsys@transformshift{3.518414in}{0.352669in}%
\pgfsys@useobject{currentmarker}{}%
\end{pgfscope}%
\end{pgfscope}%
\begin{pgfscope}%
\pgfsetbuttcap%
\pgfsetroundjoin%
\definecolor{currentfill}{rgb}{0.000000,0.000000,0.000000}%
\pgfsetfillcolor{currentfill}%
\pgfsetlinewidth{0.301125pt}%
\definecolor{currentstroke}{rgb}{0.000000,0.000000,0.000000}%
\pgfsetstrokecolor{currentstroke}%
\pgfsetdash{}{0pt}%
\pgfsys@defobject{currentmarker}{\pgfqpoint{0.000000in}{0.000000in}}{\pgfqpoint{0.000000in}{0.027778in}}{%
\pgfpathmoveto{\pgfqpoint{0.000000in}{0.000000in}}%
\pgfpathlineto{\pgfqpoint{0.000000in}{0.027778in}}%
\pgfusepath{stroke,fill}%
}%
\begin{pgfscope}%
\pgfsys@transformshift{3.613185in}{0.352669in}%
\pgfsys@useobject{currentmarker}{}%
\end{pgfscope}%
\end{pgfscope}%
\begin{pgfscope}%
\definecolor{textcolor}{rgb}{0.000000,0.000000,0.000000}%
\pgfsetstrokecolor{textcolor}%
\pgfsetfillcolor{textcolor}%
\pgftext[x=2.120801in,y=0.140972in,,top]{\color{textcolor}{\sffamily\fontsize{9.000000}{10.800000}\selectfont\catcode`\^=\active\def^{\ifmmode\sp\else\^{}\fi}\catcode`\%=\active\def%{\%}$\lambda$}}%
\end{pgfscope}%
\begin{pgfscope}%
\pgfsetbuttcap%
\pgfsetroundjoin%
\definecolor{currentfill}{rgb}{0.000000,0.000000,0.000000}%
\pgfsetfillcolor{currentfill}%
\pgfsetlinewidth{0.501875pt}%
\definecolor{currentstroke}{rgb}{0.000000,0.000000,0.000000}%
\pgfsetstrokecolor{currentstroke}%
\pgfsetdash{}{0pt}%
\pgfsys@defobject{currentmarker}{\pgfqpoint{0.000000in}{0.000000in}}{\pgfqpoint{0.041667in}{0.000000in}}{%
\pgfpathmoveto{\pgfqpoint{0.000000in}{0.000000in}}%
\pgfpathlineto{\pgfqpoint{0.041667in}{0.000000in}}%
\pgfusepath{stroke,fill}%
}%
\begin{pgfscope}%
\pgfsys@transformshift{0.530750in}{0.675077in}%
\pgfsys@useobject{currentmarker}{}%
\end{pgfscope}%
\end{pgfscope}%
\begin{pgfscope}%
\definecolor{textcolor}{rgb}{0.000000,0.000000,0.000000}%
\pgfsetstrokecolor{textcolor}%
\pgfsetfillcolor{textcolor}%
\pgftext[x=0.198444in, y=0.632868in, left, base]{\color{textcolor}{\sffamily\fontsize{8.000000}{9.600000}\selectfont\catcode`\^=\active\def^{\ifmmode\sp\else\^{}\fi}\catcode`\%=\active\def%{\%}$\mathdefault{10^{-15}}$}}%
\end{pgfscope}%
\begin{pgfscope}%
\pgfsetbuttcap%
\pgfsetroundjoin%
\definecolor{currentfill}{rgb}{0.000000,0.000000,0.000000}%
\pgfsetfillcolor{currentfill}%
\pgfsetlinewidth{0.501875pt}%
\definecolor{currentstroke}{rgb}{0.000000,0.000000,0.000000}%
\pgfsetstrokecolor{currentstroke}%
\pgfsetdash{}{0pt}%
\pgfsys@defobject{currentmarker}{\pgfqpoint{0.000000in}{0.000000in}}{\pgfqpoint{0.041667in}{0.000000in}}{%
\pgfpathmoveto{\pgfqpoint{0.000000in}{0.000000in}}%
\pgfpathlineto{\pgfqpoint{0.041667in}{0.000000in}}%
\pgfusepath{stroke,fill}%
}%
\begin{pgfscope}%
\pgfsys@transformshift{0.530750in}{1.071402in}%
\pgfsys@useobject{currentmarker}{}%
\end{pgfscope}%
\end{pgfscope}%
\begin{pgfscope}%
\definecolor{textcolor}{rgb}{0.000000,0.000000,0.000000}%
\pgfsetstrokecolor{textcolor}%
\pgfsetfillcolor{textcolor}%
\pgftext[x=0.196527in, y=1.029192in, left, base]{\color{textcolor}{\sffamily\fontsize{8.000000}{9.600000}\selectfont\catcode`\^=\active\def^{\ifmmode\sp\else\^{}\fi}\catcode`\%=\active\def%{\%}$\mathdefault{10^{-13}}$}}%
\end{pgfscope}%
\begin{pgfscope}%
\pgfsetbuttcap%
\pgfsetroundjoin%
\definecolor{currentfill}{rgb}{0.000000,0.000000,0.000000}%
\pgfsetfillcolor{currentfill}%
\pgfsetlinewidth{0.501875pt}%
\definecolor{currentstroke}{rgb}{0.000000,0.000000,0.000000}%
\pgfsetstrokecolor{currentstroke}%
\pgfsetdash{}{0pt}%
\pgfsys@defobject{currentmarker}{\pgfqpoint{0.000000in}{0.000000in}}{\pgfqpoint{0.041667in}{0.000000in}}{%
\pgfpathmoveto{\pgfqpoint{0.000000in}{0.000000in}}%
\pgfpathlineto{\pgfqpoint{0.041667in}{0.000000in}}%
\pgfusepath{stroke,fill}%
}%
\begin{pgfscope}%
\pgfsys@transformshift{0.530750in}{1.467726in}%
\pgfsys@useobject{currentmarker}{}%
\end{pgfscope}%
\end{pgfscope}%
\begin{pgfscope}%
\definecolor{textcolor}{rgb}{0.000000,0.000000,0.000000}%
\pgfsetstrokecolor{textcolor}%
\pgfsetfillcolor{textcolor}%
\pgftext[x=0.196527in, y=1.425517in, left, base]{\color{textcolor}{\sffamily\fontsize{8.000000}{9.600000}\selectfont\catcode`\^=\active\def^{\ifmmode\sp\else\^{}\fi}\catcode`\%=\active\def%{\%}$\mathdefault{10^{-11}}$}}%
\end{pgfscope}%
\begin{pgfscope}%
\pgfsetbuttcap%
\pgfsetroundjoin%
\definecolor{currentfill}{rgb}{0.000000,0.000000,0.000000}%
\pgfsetfillcolor{currentfill}%
\pgfsetlinewidth{0.501875pt}%
\definecolor{currentstroke}{rgb}{0.000000,0.000000,0.000000}%
\pgfsetstrokecolor{currentstroke}%
\pgfsetdash{}{0pt}%
\pgfsys@defobject{currentmarker}{\pgfqpoint{0.000000in}{0.000000in}}{\pgfqpoint{0.041667in}{0.000000in}}{%
\pgfpathmoveto{\pgfqpoint{0.000000in}{0.000000in}}%
\pgfpathlineto{\pgfqpoint{0.041667in}{0.000000in}}%
\pgfusepath{stroke,fill}%
}%
\begin{pgfscope}%
\pgfsys@transformshift{0.530750in}{1.864050in}%
\pgfsys@useobject{currentmarker}{}%
\end{pgfscope}%
\end{pgfscope}%
\begin{pgfscope}%
\definecolor{textcolor}{rgb}{0.000000,0.000000,0.000000}%
\pgfsetstrokecolor{textcolor}%
\pgfsetfillcolor{textcolor}%
\pgftext[x=0.241110in, y=1.821841in, left, base]{\color{textcolor}{\sffamily\fontsize{8.000000}{9.600000}\selectfont\catcode`\^=\active\def^{\ifmmode\sp\else\^{}\fi}\catcode`\%=\active\def%{\%}$\mathdefault{10^{-9}}$}}%
\end{pgfscope}%
\begin{pgfscope}%
\definecolor{textcolor}{rgb}{0.000000,0.000000,0.000000}%
\pgfsetstrokecolor{textcolor}%
\pgfsetfillcolor{textcolor}%
\pgftext[x=0.140972in,y=1.162184in,,bottom,rotate=90.000000]{\color{textcolor}{\sffamily\fontsize{9.000000}{10.800000}\selectfont\catcode`\^=\active\def^{\ifmmode\sp\else\^{}\fi}\catcode`\%=\active\def%{\%}$|\mathcal{H}|$}}%
\end{pgfscope}%
\begin{pgfscope}%
\pgfpathrectangle{\pgfqpoint{0.530750in}{0.352669in}}{\pgfqpoint{3.180103in}{1.619031in}}%
\pgfusepath{clip}%
\pgfsetrectcap%
\pgfsetroundjoin%
\pgfsetlinewidth{0.803000pt}%
\definecolor{currentstroke}{rgb}{0.768627,0.556863,0.211765}%
\pgfsetstrokecolor{currentstroke}%
\pgfsetdash{}{0pt}%
\pgfpathmoveto{\pgfqpoint{0.675300in}{0.485914in}}%
\pgfpathlineto{\pgfqpoint{0.675392in}{0.426261in}}%
\pgfpathlineto{\pgfqpoint{0.676319in}{0.520808in}}%
\pgfpathlineto{\pgfqpoint{0.685584in}{0.485914in}}%
\pgfpathlineto{\pgfqpoint{0.749333in}{1.028715in}}%
\pgfpathlineto{\pgfqpoint{0.815638in}{1.137587in}}%
\pgfpathlineto{\pgfqpoint{0.878376in}{1.204609in}}%
\pgfpathlineto{\pgfqpoint{0.936272in}{1.254496in}}%
\pgfpathlineto{\pgfqpoint{0.989104in}{1.294807in}}%
\pgfpathlineto{\pgfqpoint{1.037065in}{1.328808in}}%
\pgfpathlineto{\pgfqpoint{1.119906in}{1.384367in}}%
\pgfpathlineto{\pgfqpoint{1.314037in}{1.511983in}}%
\pgfpathlineto{\pgfqpoint{1.414262in}{1.578914in}}%
\pgfpathlineto{\pgfqpoint{1.440181in}{1.594016in}}%
\pgfpathlineto{\pgfqpoint{1.452064in}{1.599900in}}%
\pgfpathlineto{\pgfqpoint{1.463281in}{1.604568in}}%
\pgfpathlineto{\pgfqpoint{1.473866in}{1.607994in}}%
\pgfpathlineto{\pgfqpoint{1.483853in}{1.610228in}}%
\pgfpathlineto{\pgfqpoint{1.493282in}{1.611377in}}%
\pgfpathlineto{\pgfqpoint{1.502201in}{1.611590in}}%
\pgfpathlineto{\pgfqpoint{1.518711in}{1.609895in}}%
\pgfpathlineto{\pgfqpoint{1.554474in}{1.602342in}}%
\pgfpathlineto{\pgfqpoint{1.560997in}{1.602268in}}%
\pgfpathlineto{\pgfqpoint{1.567374in}{1.603417in}}%
\pgfpathlineto{\pgfqpoint{1.573618in}{1.606051in}}%
\pgfpathlineto{\pgfqpoint{1.579734in}{1.610249in}}%
\pgfpathlineto{\pgfqpoint{1.585718in}{1.615839in}}%
\pgfpathlineto{\pgfqpoint{1.597274in}{1.629525in}}%
\pgfpathlineto{\pgfqpoint{1.613626in}{1.650111in}}%
\pgfpathlineto{\pgfqpoint{1.623504in}{1.658473in}}%
\pgfpathlineto{\pgfqpoint{1.632743in}{1.663765in}}%
\pgfpathlineto{\pgfqpoint{1.641574in}{1.667221in}}%
\pgfpathlineto{\pgfqpoint{1.654386in}{1.670207in}}%
\pgfpathlineto{\pgfqpoint{1.667006in}{1.671477in}}%
\pgfpathlineto{\pgfqpoint{1.684063in}{1.671614in}}%
\pgfpathlineto{\pgfqpoint{1.725882in}{1.671033in}}%
\pgfpathlineto{\pgfqpoint{1.803247in}{1.673135in}}%
\pgfpathlineto{\pgfqpoint{1.819671in}{1.671155in}}%
\pgfpathlineto{\pgfqpoint{1.847578in}{1.664291in}}%
\pgfpathlineto{\pgfqpoint{1.871404in}{1.662354in}}%
\pgfpathlineto{\pgfqpoint{1.936378in}{1.657071in}}%
\pgfpathlineto{\pgfqpoint{2.040992in}{1.648599in}}%
\pgfpathlineto{\pgfqpoint{2.112028in}{1.644534in}}%
\pgfpathlineto{\pgfqpoint{2.176760in}{1.642348in}}%
\pgfpathlineto{\pgfqpoint{2.253693in}{1.641429in}}%
\pgfpathlineto{\pgfqpoint{2.345365in}{1.642080in}}%
\pgfpathlineto{\pgfqpoint{2.454833in}{1.644305in}}%
\pgfpathlineto{\pgfqpoint{2.743483in}{1.652377in}}%
\pgfpathlineto{\pgfqpoint{3.082957in}{1.661164in}}%
\pgfpathlineto{\pgfqpoint{3.253242in}{1.664900in}}%
\pgfpathlineto{\pgfqpoint{3.253242in}{1.664900in}}%
\pgfusepath{stroke}%
\end{pgfscope}%
\begin{pgfscope}%
\pgfpathrectangle{\pgfqpoint{0.530750in}{0.352669in}}{\pgfqpoint{3.180103in}{1.619031in}}%
\pgfusepath{clip}%
\pgfsetrectcap%
\pgfsetroundjoin%
\pgfsetlinewidth{0.803000pt}%
\definecolor{currentstroke}{rgb}{0.250980,0.400000,0.600000}%
\pgfsetstrokecolor{currentstroke}%
\pgfsetdash{}{0pt}%
\pgfpathmoveto{\pgfqpoint{0.675392in}{0.346002in}}%
\pgfpathlineto{\pgfqpoint{0.675393in}{0.426261in}}%
\pgfpathlineto{\pgfqpoint{0.685600in}{0.426261in}}%
\pgfpathlineto{\pgfqpoint{0.764822in}{0.984334in}}%
\pgfpathlineto{\pgfqpoint{0.847157in}{1.101746in}}%
\pgfpathlineto{\pgfqpoint{0.913590in}{1.146836in}}%
\pgfpathlineto{\pgfqpoint{1.025832in}{1.218019in}}%
\pgfpathlineto{\pgfqpoint{1.115296in}{1.274162in}}%
\pgfpathlineto{\pgfqpoint{1.187235in}{1.321071in}}%
\pgfpathlineto{\pgfqpoint{1.221297in}{1.355108in}}%
\pgfpathlineto{\pgfqpoint{1.248779in}{1.372332in}}%
\pgfpathlineto{\pgfqpoint{1.323776in}{1.444543in}}%
\pgfpathlineto{\pgfqpoint{1.434841in}{1.548630in}}%
\pgfpathlineto{\pgfqpoint{1.492686in}{1.604006in}}%
\pgfpathlineto{\pgfqpoint{1.514357in}{1.622316in}}%
\pgfpathlineto{\pgfqpoint{1.526785in}{1.630969in}}%
\pgfpathlineto{\pgfqpoint{1.537827in}{1.637126in}}%
\pgfpathlineto{\pgfqpoint{1.552191in}{1.643012in}}%
\pgfpathlineto{\pgfqpoint{1.584157in}{1.653863in}}%
\pgfpathlineto{\pgfqpoint{1.594769in}{1.659505in}}%
\pgfpathlineto{\pgfqpoint{1.613593in}{1.671654in}}%
\pgfpathlineto{\pgfqpoint{1.622479in}{1.679410in}}%
\pgfpathlineto{\pgfqpoint{1.632286in}{1.689950in}}%
\pgfpathlineto{\pgfqpoint{1.640060in}{1.700650in}}%
\pgfpathlineto{\pgfqpoint{1.649069in}{1.715846in}}%
\pgfpathlineto{\pgfqpoint{1.663067in}{1.739826in}}%
\pgfpathlineto{\pgfqpoint{1.669371in}{1.747810in}}%
\pgfpathlineto{\pgfqpoint{1.675173in}{1.753187in}}%
\pgfpathlineto{\pgfqpoint{1.694386in}{1.768239in}}%
\pgfpathlineto{\pgfqpoint{1.704369in}{1.780012in}}%
\pgfpathlineto{\pgfqpoint{1.719855in}{1.798789in}}%
\pgfpathlineto{\pgfqpoint{1.727068in}{1.804653in}}%
\pgfpathlineto{\pgfqpoint{1.733734in}{1.808160in}}%
\pgfpathlineto{\pgfqpoint{1.742105in}{1.812086in}}%
\pgfpathlineto{\pgfqpoint{1.750063in}{1.817535in}}%
\pgfpathlineto{\pgfqpoint{1.760497in}{1.826574in}}%
\pgfpathlineto{\pgfqpoint{1.771003in}{1.835415in}}%
\pgfpathlineto{\pgfqpoint{1.779000in}{1.840211in}}%
\pgfpathlineto{\pgfqpoint{1.786401in}{1.842987in}}%
\pgfpathlineto{\pgfqpoint{1.794605in}{1.845749in}}%
\pgfpathlineto{\pgfqpoint{1.804476in}{1.850758in}}%
\pgfpathlineto{\pgfqpoint{1.830217in}{1.865478in}}%
\pgfpathlineto{\pgfqpoint{1.838585in}{1.867837in}}%
\pgfpathlineto{\pgfqpoint{1.848785in}{1.870510in}}%
\pgfpathlineto{\pgfqpoint{1.859961in}{1.874990in}}%
\pgfpathlineto{\pgfqpoint{1.880109in}{1.883566in}}%
\pgfpathlineto{\pgfqpoint{1.889992in}{1.885839in}}%
\pgfpathlineto{\pgfqpoint{1.907903in}{1.889158in}}%
\pgfpathlineto{\pgfqpoint{1.939720in}{1.896305in}}%
\pgfpathlineto{\pgfqpoint{1.954852in}{1.897299in}}%
\pgfpathlineto{\pgfqpoint{1.991776in}{1.898080in}}%
\pgfpathlineto{\pgfqpoint{2.026126in}{1.897011in}}%
\pgfpathlineto{\pgfqpoint{2.158598in}{1.891478in}}%
\pgfpathlineto{\pgfqpoint{2.244295in}{1.890393in}}%
\pgfpathlineto{\pgfqpoint{2.393865in}{1.890150in}}%
\pgfpathlineto{\pgfqpoint{3.129481in}{1.891710in}}%
\pgfpathlineto{\pgfqpoint{3.566302in}{1.892431in}}%
\pgfpathlineto{\pgfqpoint{3.566302in}{1.892431in}}%
\pgfusepath{stroke}%
\end{pgfscope}%
\begin{pgfscope}%
\pgfpathrectangle{\pgfqpoint{0.530750in}{0.352669in}}{\pgfqpoint{3.180103in}{1.619031in}}%
\pgfusepath{clip}%
\pgfsetrectcap%
\pgfsetroundjoin%
\pgfsetlinewidth{0.803000pt}%
\definecolor{currentstroke}{rgb}{0.619608,0.188235,0.172549}%
\pgfsetstrokecolor{currentstroke}%
\pgfsetdash{}{0pt}%
\pgfpathmoveto{\pgfqpoint{0.675300in}{0.426261in}}%
\pgfpathlineto{\pgfqpoint{0.675393in}{0.426261in}}%
\pgfpathlineto{\pgfqpoint{0.675393in}{0.346002in}}%
\pgfpathmoveto{\pgfqpoint{0.685600in}{0.346002in}}%
\pgfpathlineto{\pgfqpoint{0.685604in}{0.426261in}}%
\pgfpathlineto{\pgfqpoint{0.775450in}{0.949110in}}%
\pgfpathlineto{\pgfqpoint{0.858818in}{1.039468in}}%
\pgfpathlineto{\pgfqpoint{0.931944in}{1.097018in}}%
\pgfpathlineto{\pgfqpoint{0.995958in}{1.141778in}}%
\pgfpathlineto{\pgfqpoint{1.052055in}{1.179378in}}%
\pgfpathlineto{\pgfqpoint{1.101341in}{1.212163in}}%
\pgfpathlineto{\pgfqpoint{1.144792in}{1.241436in}}%
\pgfpathlineto{\pgfqpoint{1.183244in}{1.267966in}}%
\pgfpathlineto{\pgfqpoint{1.217409in}{1.292236in}}%
\pgfpathlineto{\pgfqpoint{1.247887in}{1.314642in}}%
\pgfpathlineto{\pgfqpoint{1.275185in}{1.335472in}}%
\pgfpathlineto{\pgfqpoint{1.299731in}{1.354920in}}%
\pgfpathlineto{\pgfqpoint{1.324277in}{1.384919in}}%
\pgfpathlineto{\pgfqpoint{1.344114in}{1.400243in}}%
\pgfpathlineto{\pgfqpoint{1.363951in}{1.424062in}}%
\pgfpathlineto{\pgfqpoint{1.382007in}{1.445961in}}%
\pgfpathlineto{\pgfqpoint{1.398322in}{1.465293in}}%
\pgfpathlineto{\pgfqpoint{1.413122in}{1.482616in}}%
\pgfpathlineto{\pgfqpoint{1.426603in}{1.498355in}}%
\pgfpathlineto{\pgfqpoint{1.438933in}{1.512796in}}%
\pgfpathlineto{\pgfqpoint{1.450252in}{1.526149in}}%
\pgfpathlineto{\pgfqpoint{1.460678in}{1.538562in}}%
\pgfpathlineto{\pgfqpoint{1.470314in}{1.550149in}}%
\pgfpathlineto{\pgfqpoint{1.479246in}{1.560995in}}%
\pgfpathlineto{\pgfqpoint{1.487548in}{1.571161in}}%
\pgfpathlineto{\pgfqpoint{1.495284in}{1.580693in}}%
\pgfpathlineto{\pgfqpoint{1.502508in}{1.589623in}}%
\pgfpathlineto{\pgfqpoint{1.509269in}{1.597973in}}%
\pgfpathlineto{\pgfqpoint{1.515605in}{1.605755in}}%
\pgfpathlineto{\pgfqpoint{1.521552in}{1.612978in}}%
\pgfpathlineto{\pgfqpoint{1.527139in}{1.619646in}}%
\pgfpathlineto{\pgfqpoint{1.532393in}{1.625767in}}%
\pgfpathlineto{\pgfqpoint{1.537335in}{1.631353in}}%
\pgfpathlineto{\pgfqpoint{1.541987in}{1.636420in}}%
\pgfpathlineto{\pgfqpoint{1.546365in}{1.640996in}}%
\pgfpathlineto{\pgfqpoint{1.550488in}{1.645115in}}%
\pgfpathlineto{\pgfqpoint{1.554372in}{1.648819in}}%
\pgfpathlineto{\pgfqpoint{1.558031in}{1.652156in}}%
\pgfpathlineto{\pgfqpoint{1.561480in}{1.655175in}}%
\pgfpathlineto{\pgfqpoint{1.564734in}{1.657923in}}%
\pgfpathlineto{\pgfqpoint{1.567806in}{1.660444in}}%
\pgfpathlineto{\pgfqpoint{1.570708in}{1.662779in}}%
\pgfpathlineto{\pgfqpoint{1.573453in}{1.664961in}}%
\pgfpathlineto{\pgfqpoint{1.576051in}{1.667020in}}%
\pgfpathlineto{\pgfqpoint{1.578514in}{1.668978in}}%
\pgfpathlineto{\pgfqpoint{1.580850in}{1.670853in}}%
\pgfpathlineto{\pgfqpoint{1.583069in}{1.672660in}}%
\pgfpathlineto{\pgfqpoint{1.585179in}{1.674407in}}%
\pgfpathlineto{\pgfqpoint{1.587187in}{1.676103in}}%
\pgfpathlineto{\pgfqpoint{1.589100in}{1.677752in}}%
\pgfpathlineto{\pgfqpoint{1.590926in}{1.679358in}}%
\pgfpathlineto{\pgfqpoint{1.592669in}{1.680923in}}%
\pgfpathlineto{\pgfqpoint{1.594335in}{1.682449in}}%
\pgfpathlineto{\pgfqpoint{1.595928in}{1.683935in}}%
\pgfpathlineto{\pgfqpoint{1.597454in}{1.685381in}}%
\pgfpathlineto{\pgfqpoint{1.598917in}{1.686788in}}%
\pgfpathlineto{\pgfqpoint{1.600319in}{1.688155in}}%
\pgfpathlineto{\pgfqpoint{1.601664in}{1.689481in}}%
\pgfpathlineto{\pgfqpoint{1.602956in}{1.690764in}}%
\pgfpathlineto{\pgfqpoint{1.604198in}{1.692005in}}%
\pgfpathlineto{\pgfqpoint{1.605390in}{1.693201in}}%
\pgfpathlineto{\pgfqpoint{1.606537in}{1.694351in}}%
\pgfpathlineto{\pgfqpoint{1.607640in}{1.695455in}}%
\pgfpathlineto{\pgfqpoint{1.608708in}{1.696561in}}%
\pgfpathlineto{\pgfqpoint{1.609752in}{1.697721in}}%
\pgfpathlineto{\pgfqpoint{1.610769in}{1.698928in}}%
\pgfpathlineto{\pgfqpoint{1.611761in}{1.700172in}}%
\pgfpathlineto{\pgfqpoint{1.612726in}{1.701444in}}%
\pgfpathlineto{\pgfqpoint{1.613664in}{1.702736in}}%
\pgfpathlineto{\pgfqpoint{1.614577in}{1.704043in}}%
\pgfpathlineto{\pgfqpoint{1.615463in}{1.705358in}}%
\pgfpathlineto{\pgfqpoint{1.616327in}{1.706704in}}%
\pgfpathlineto{\pgfqpoint{1.617171in}{1.708096in}}%
\pgfpathlineto{\pgfqpoint{1.617995in}{1.709535in}}%
\pgfpathlineto{\pgfqpoint{1.618799in}{1.711023in}}%
\pgfpathlineto{\pgfqpoint{1.619585in}{1.712562in}}%
\pgfpathlineto{\pgfqpoint{1.620352in}{1.714155in}}%
\pgfpathlineto{\pgfqpoint{1.621102in}{1.715804in}}%
\pgfusepath{stroke}%
\end{pgfscope}%
\begin{pgfscope}%
\pgfsetrectcap%
\pgfsetmiterjoin%
\pgfsetlinewidth{0.602250pt}%
\definecolor{currentstroke}{rgb}{0.000000,0.000000,0.000000}%
\pgfsetstrokecolor{currentstroke}%
\pgfsetdash{}{0pt}%
\pgfpathmoveto{\pgfqpoint{0.530750in}{0.352669in}}%
\pgfpathlineto{\pgfqpoint{0.530750in}{1.971700in}}%
\pgfusepath{stroke}%
\end{pgfscope}%
\begin{pgfscope}%
\pgfsetrectcap%
\pgfsetmiterjoin%
\pgfsetlinewidth{0.602250pt}%
\definecolor{currentstroke}{rgb}{0.000000,0.000000,0.000000}%
\pgfsetstrokecolor{currentstroke}%
\pgfsetdash{}{0pt}%
\pgfpathmoveto{\pgfqpoint{0.530750in}{0.352669in}}%
\pgfpathlineto{\pgfqpoint{3.710853in}{0.352669in}}%
\pgfusepath{stroke}%
\end{pgfscope}%
\begin{pgfscope}%
\definecolor{textcolor}{rgb}{0.431373,0.423529,0.462745}%
\pgfsetstrokecolor{textcolor}%
\pgfsetfillcolor{textcolor}%
\pgftext[x=3.508482in,y=1.706337in,right,bottom]{\color{textcolor}{\sffamily\fontsize{6.000000}{7.200000}\selectfont\catcode`\^=\active\def^{\ifmmode\sp\else\^{}\fi}\catcode`\%=\active\def%{\%}$\epsilon = 10^{-10}$}}%
\end{pgfscope}%
\begin{pgfscope}%
\pgfsetrectcap%
\pgfsetroundjoin%
\pgfsetlinewidth{0.803000pt}%
\definecolor{currentstroke}{rgb}{0.768627,0.556863,0.211765}%
\pgfsetstrokecolor{currentstroke}%
\pgfsetdash{}{0pt}%
\pgfpathmoveto{\pgfqpoint{3.065589in}{0.741409in}}%
\pgfpathlineto{\pgfqpoint{3.155867in}{0.741409in}}%
\pgfpathlineto{\pgfqpoint{3.246144in}{0.741409in}}%
\pgfusepath{stroke}%
\end{pgfscope}%
\begin{pgfscope}%
\definecolor{textcolor}{rgb}{0.000000,0.000000,0.000000}%
\pgfsetstrokecolor{textcolor}%
\pgfsetfillcolor{textcolor}%
\pgftext[x=3.318367in,y=0.709812in,left,base]{\color{textcolor}{\sffamily\fontsize{6.500000}{7.800000}\selectfont\catcode`\^=\active\def^{\ifmmode\sp\else\^{}\fi}\catcode`\%=\active\def%{\%}$b = 8M$}}%
\end{pgfscope}%
\begin{pgfscope}%
\pgfsetrectcap%
\pgfsetroundjoin%
\pgfsetlinewidth{0.803000pt}%
\definecolor{currentstroke}{rgb}{0.250980,0.400000,0.600000}%
\pgfsetstrokecolor{currentstroke}%
\pgfsetdash{}{0pt}%
\pgfpathmoveto{\pgfqpoint{3.065589in}{0.608902in}}%
\pgfpathlineto{\pgfqpoint{3.155867in}{0.608902in}}%
\pgfpathlineto{\pgfqpoint{3.246144in}{0.608902in}}%
\pgfusepath{stroke}%
\end{pgfscope}%
\begin{pgfscope}%
\definecolor{textcolor}{rgb}{0.000000,0.000000,0.000000}%
\pgfsetstrokecolor{textcolor}%
\pgfsetfillcolor{textcolor}%
\pgftext[x=3.318367in,y=0.577305in,left,base]{\color{textcolor}{\sffamily\fontsize{6.500000}{7.800000}\selectfont\catcode`\^=\active\def^{\ifmmode\sp\else\^{}\fi}\catcode`\%=\active\def%{\%}$b \approx b_c$}}%
\end{pgfscope}%
\begin{pgfscope}%
\pgfsetrectcap%
\pgfsetroundjoin%
\pgfsetlinewidth{0.803000pt}%
\definecolor{currentstroke}{rgb}{0.619608,0.188235,0.172549}%
\pgfsetstrokecolor{currentstroke}%
\pgfsetdash{}{0pt}%
\pgfpathmoveto{\pgfqpoint{3.065589in}{0.475267in}}%
\pgfpathlineto{\pgfqpoint{3.155867in}{0.475267in}}%
\pgfpathlineto{\pgfqpoint{3.246144in}{0.475267in}}%
\pgfusepath{stroke}%
\end{pgfscope}%
\begin{pgfscope}%
\definecolor{textcolor}{rgb}{0.000000,0.000000,0.000000}%
\pgfsetstrokecolor{textcolor}%
\pgfsetfillcolor{textcolor}%
\pgftext[x=3.318367in,y=0.443669in,left,base]{\color{textcolor}{\sffamily\fontsize{6.500000}{7.800000}\selectfont\catcode`\^=\active\def^{\ifmmode\sp\else\^{}\fi}\catcode`\%=\active\def%{\%}$b = 4M$}}%
\end{pgfscope}%
\end{pgfpicture}%
\makeatother%
\endgroup%

  \caption[Hamiltonian constraint violation for three null geodesics]{%
    Hamiltonian constraint violation $\abs{\mathcal{H}}$ as a function of
    affine parameter $\lambda$ for three equatorial null geodesics in
    Schwarzschild spacetime ($M = 1$), integrated with RK45 at tolerances
    $\epsilon_a = \epsilon_r = 10^{-10}$.  Gold: the deflected ray
    ($b = 8M$) showing steady, gradual growth that saturates near the
    tolerance (dashed line).  Blue: the near-critical ray
    ($b \approx b_c$) whose violation grows faster during the
    photon-sphere passage ($\lambda \approx 45$--$55$) as the Lyapunov
    instability amplifies local errors, reaching roughly $10\,\epsilon_a$.
    Crimson: the plunge ray ($b = 4M$) whose constraint violation
    steepens as it approaches the horizon.  Observer at $r_0 = 50M$.}
  \label{fig:H-violation}
\end{figure}

The practical consequence is that conservation monitoring is most
valuable not as a per-ray abort criterion but as a global quality
metric: the worst-case $\abs{\mathcal{H}}$ across an entire image
reveals whether the tolerances were adequate for that particular
geometry and observer configuration.  The codebase sets tolerances
conservatively enough that all rays satisfy the bound
\cref{eq:H-violation-bound} by construction; after rendering, the
maximum $\abs{\mathcal{H}}$ across all rays confirms that no
trajectory exceeded the error budget.
\xrefnote{Convergence studies verifying $\abs{\mathcal{H}}$ against
tolerance using \texttt{traceGeodesic} are developed in
\cref{ch:ode-integration}.}
